\title{2024 年高考数学 北京卷}
\date{}
\maketitle

\section*{一、单项选择题}
本部分共10小题，每小题4分，共40分。在每小题列出的四个选项中，选出符合题目要求的一项。

\begin{question}[points=4]
已知集合 $M = \{ x \mid -4 < x \leq 1 \}$，$N = \{x \mid -1 < x < 3\}$，则 $M \cup N =$ （ ）
\begin{choices}
  \item $\{x \mid -4 < x < 3\}$
  \item $\{x \mid -1 < x \leq 1\}$
  \item $\{0,1,2\}$
  \item $\{x \mid -1 < x < 4\}$
\end{choices}
\begin{solution}
\textbf{答案：}A

\textbf{分析：}直接根据并集含义即可得到答案。

\textbf{详解：}由题意得 $M \cup N = (-4,3)$，故选：A。
\end{solution}
\end{question}

\begin{question}[points=4]
已知 $\dfrac{z}{i} = i-1$，则 $z=$ （ ）
\begin{choices}
  \item $1-i$
  \item $-i$
  \item $-1-i$
  \item $1$
\end{choices}
\begin{solution}
\textbf{答案：}C

\textbf{分析：}直接根据复数乘法即可得到答案。

\textbf{详解：}由题意得 $z = i(i-1) = -1-i$，故选：C。
\end{solution}
\end{question}

\begin{question}[points=4]
求圆 $x^2 + y^2 - 2x + 6y = 0$ 的圆心到 $x - y + 2 = 0$ 的距离（ ）
\begin{choices}
  \item $2\sqrt{3}$
  \item $2$
  \item $3\sqrt{2}$
  \item $\sqrt{6}$
\end{choices}
\begin{solution}
\textbf{答案：}C

\textbf{分析：}求出圆心坐标，再利用点到直线距离公式即可。

\textbf{详解：}由题意得 $x^2 + y^2 - 2x + 6y = 0$，即 $(x-1)^2 + (y+3)^2 = 10$，则其圆心坐标为 $(1,-3)$，则圆心到直线 $x-y+2=0$ 的距离为 $\dfrac{|1+3+2|}{\sqrt{1^2+1^2}} = 3\sqrt{2}$，故选：C。
\end{solution}
\end{question}

\begin{question}[points=4]
$(x - \sqrt{x})^4$ 的二项展开式中 $x^3$ 的系数为（ ）
\begin{choices}
  \item $15$
  \item $6$
  \item $-4$
  \item $-13$
\end{choices}
\begin{solution}
\textbf{答案：}B

\textbf{分析：}写出二项展开式，令 $4-\frac{r}{2}=3$，解出 $r$ 然后回代入二项展开式系数即可得解。

\textbf{详解：}$(x - \sqrt{x})^4$ 的二项展开式为 $T_{r+1} = C_4^r x^{4-r} (-\sqrt{x})^r = C_4^r (-1)^r x^{4-\frac{r}{2}}$，$r=0,1,2,3,4$。令 $4-\frac{r}{2}=3$，解得 $r=2$，故所求即为 $C_4^2 (-1)^2 = 6$。故选：B。
\end{solution}
\end{question}

\begin{question}[points=4]
已知向量 $\vec{a}$，$\vec{b}$，则“$(\vec{a}+\vec{b})\cdot(\vec{a}-\vec{b})=0$”是“$\vec{a}=\vec{b}$ 或 $\vec{a}=-\vec{b}$”的（ ）条件。
\begin{choices}
  \item 必要而不充分条件
  \item 充分而不必要条件
  \item 充分且必要条件
  \item 既不充分也不必要条件
\end{choices}
\begin{solution}
\textbf{答案：}A

\textbf{分析：}根据向量数量积分析可知 $(\vec{a}+\vec{b})\cdot(\vec{a}-\vec{b})=0$ 等价于 $|\vec{a}|=|\vec{b}|$，结合充分、必要条件分析判断。

\textbf{详解：}因为 $(\vec{a}+\vec{b})\cdot(\vec{a}-\vec{b}) = |\vec{a}|^2-|\vec{b}|^2=0$，可得 $|\vec{a}|^2=|\vec{b}|^2$，即 $|\vec{a}|=|\vec{b}|$，可知 $(\vec{a}+\vec{b})\cdot(\vec{a}-\vec{b})=0$ 等价于 $|\vec{a}|=|\vec{b}|$。若 $\vec{a}=\vec{b}$ 或 $\vec{a}=-\vec{b}$，可得 $|\vec{a}|=|\vec{b}|$，即 $(\vec{a}+\vec{b})\cdot(\vec{a}-\vec{b})=0$，可知必要性成立；若 $(\vec{a}+\vec{b})\cdot(\vec{a}-\vec{b})=0$，即 $|\vec{a}|=|\vec{b}|$，无法得出 $\vec{a}=\vec{b}$ 或 $\vec{a}=-\vec{b}$，例如 $\vec{a}=(1,0),\vec{b}=(0,1)$，满足 $|\vec{a}|=|\vec{b}|$，但 $\vec{a}\neq\vec{b}$ 且 $\vec{a}\neq-\vec{b}$，可知充分性不成立。综上所述，“$(\vec{a}+\vec{b})\cdot(\vec{a}-\vec{b})=0$”是“$\vec{a}=\vec{b}$ 或 $\vec{a}=-\vec{b}$”的必要不充分条件。故选：A。
\end{solution}
\end{question}

\begin{question}[points=4]
已知 $f(x)=\sin\omega x(\omega>0)$，$f(x_1)=-1$，$f(x_2)=1$，$|x_1-x_2|_{\min}=\dfrac{\pi}{2}$，则 $\omega=$ （ ）
\begin{choices}
  \item $1$
  \item $2$
  \item $3$
  \item $4$
\end{choices}
\begin{solution}
\textbf{答案：}B

\textbf{分析：}根据三角函数最值分析周期性，结合三角函数最小正周期公式运算求解。

\textbf{详解：}由题意可知：$x_1$ 为 $f(x)$ 的最小值点，$x_2$ 为 $f(x)$ 的最大值点，则 $|x_1-x_2|_{\min}=\dfrac{T}{2}=\dfrac{\pi}{2}$，即 $T=\pi$，且 $\omega>0$，所以 $\omega=\dfrac{2\pi}{T}=2$。故选：B。
\end{solution}
\end{question}

\begin{question}[points=4]
记水的质量为 $d=\dfrac{S-1}{\ln n}$，并且 $d$ 越大，水质量越好。若 $S$ 不变，且 $d_1=2.1$，$d_2=2.2$，则 $n_1$ 与 $n_2$ 的关系为（ ）
\begin{choices}
  \item $n_1 < n_2$
  \item $n_1 > n_2$
  \item 若 $S<1$，则 $n_1<n_2$；若 $S>1$，则 $n_1>n_2$；
  \item 若 $S<1$，则 $n_1>n_2$；若 $S>1$，则 $n_1<n_2$；
\end{choices}
\begin{solution}
\textbf{答案：}C

\textbf{分析：}根据题意分析可得 $\begin{cases} n_1 = e^{\frac{S-1}{2.1}} \\ n_2 = e^{\frac{S-1}{2.2}} \end{cases}$，讨论 $S$ 与 $1$ 的大小关系，结合指数函数单调性分析判断。

\textbf{详解：}由题意可得 $\begin{cases} d_1 = \dfrac{S-1}{\ln n_1}=2.1 \\ d_2 = \dfrac{S-1}{\ln n_2}=2.2 \end{cases}$，解得 $\begin{cases} n_1 = e^{\frac{S-1}{2.1}} \\ n_2 = e^{\frac{S-1}{2.2}} \end{cases}$。若 $S>1$，则 $\dfrac{S-1}{2.1} > \dfrac{S-1}{2.2}$，可得 $e^{\frac{S-1}{2.1}} > e^{\frac{S-1}{2.2}}$，即 $n_1>n_2$；若 $S=1$，则 $\dfrac{S-1}{2.1} = \dfrac{S-1}{2.2} =0$，可得 $n_1=n_2=1$；若 $S<1$，则 $\dfrac{S-1}{2.1} < \dfrac{S-1}{2.2}$，可得 $e^{\frac{S-1}{2.1}} < e^{\frac{S-1}{2.2}}$，即 $n_1<n_2$。结合选项可知 C 正确，ABD 错误。故选：C。
\end{solution}
\end{question}

\begin{question}[points=4]
已知以边长为4的正方形为底面的四棱锥，四条侧棱分别为4，4，$2\sqrt{2}$，$2\sqrt{2}$，则该四棱锥的高为（ ）
\begin{choices}
  \item $\dfrac{\sqrt{2}}{2}$
  \item $\dfrac{\sqrt{3}}{2}$
  \item $2\sqrt{3}$
  \item $\sqrt{3}$
\end{choices}
\begin{solution}
\textbf{答案：}D

\textbf{分析：}取点作辅助线，根据题意分析可知平面 $PEF \perp$ 平面 $ABCD$，可知 $PO \perp$ 平面 $ABCD$，利用等体积法求点到面的距离。

\textbf{详解：}如图，底面 $ABCD$ 为正方形，当相邻的棱长相等时，不妨设 $PA=PB=AB=4, PC=PD=2\sqrt{2}$。分别取 $AB, CD$ 的中点 $E,F$，连接 $PE, PF, EF$，则 $PE \perp AB, EF \perp AB$，且 $PE \cap EF = E$，$PE, EF \subset$ 平面 $PEF$，可知 $AB \perp$ 平面 $PEF$，且 $AB \subset$ 平面 $ABCD$，所以平面 $PEF \perp$ 平面 $ABCD$。过 $P$ 作 $EF$ 的垂线，垂足为 $O$，即 $PO \perp EF$，由平面 $PEF \cap$ 平面 $ABCD = EF$，$PO \subset$ 平面 $PEF$，所以 $PO \perp$ 平面 $ABCD$。由题意可得：$PE=2\sqrt{3}, PF=2, EF=4$，则 $PE^2+PF^2=EF^2$，即 $PE \perp PF$，则 $\dfrac{1}{2} PE \cdot PF = \dfrac{1}{2} PO \cdot EF$，可得 $PO = \dfrac{PE \cdot PF}{EF} = \sqrt{3}$。所以四棱锥的高为 $\sqrt{3}$。当相对的棱长相等时，不妨设 $PA=PC=4, PB=PD=2\sqrt{2}$，因为 $BD=4\sqrt{2}=PB+PD$，此时不能形成三角形 $PBD$，与题意不符，这种情况不存在。故选：D。
\end{solution}
\end{question}

\begin{question}[points=4]
已知 $(x_1,y_1)$，$(x_2,y_2)$ 是函数 $y=2^x$ 图象上不同的两点，则下列正确的是（ ）
\begin{choices}
  \item $\log_2 \dfrac{y_1+y_2}{2} > \dfrac{x_1+x_2}{2}$
  \item $\log_2 \dfrac{y_1+y_2}{2} < \dfrac{x_1+x_2}{2}$
  \item $\log_2 \dfrac{y_1+y_2}{2} > x_1+x_2$
  \item $\log_2 \dfrac{y_1+y_2}{2} < x_1+x_2$
\end{choices}
\begin{solution}
\textbf{答案：}A

\textbf{分析：}根据指数函数和对数函数的单调性结合基本不等式分析判断AB；举例判断CD即可。

\textbf{详解：}由题意不妨设 $x_1<x_2$，因为函数 $y=2^x$ 是增函数，所以 $0<2^{x_1}<2^{x_2}$，即 $0<y_1<y_2$。对于选项AB：可得 $\dfrac{2^{x_1}+2^{x_2}}{2} > \sqrt{2^{x_1}\cdot2^{x_2}} = 2^{\frac{x_1+x_2}{2}}$，即 $\dfrac{y_1+y_2}{2} > 2^{\frac{x_1+x_2}{2}} > 0$，根据函数 $y=\log_2 x$ 是增函数，所以 $\log_2 \dfrac{y_1+y_2}{2} > \log_2 2^{\frac{x_1+x_2}{2}} = \dfrac{x_1+x_2}{2}$，故A正确，B错误。对于选项C：例如 $x_1=0, x_2=1$，则 $y_1=1, y_2=2$，可得 $\log_2 \dfrac{y_1+y_2}{2} = \log_2 \dfrac{3}{2} \in (0,1)$，即 $\log_2 \dfrac{y_1+y_2}{2} < 1 = x_1+x_2$，故C错误。对于选项D：例如 $x_1=-1, x_2=-2$，则 $y_1=\dfrac{1}{2}, y_2=\dfrac{1}{4}$，可得 $\log_2 \dfrac{y_1+y_2}{2} = \log_2 \dfrac{3}{8} = \log_2 3 - 3 \in (-2,-1)$，即 $\log_2 \dfrac{y_1+y_2}{2} > -3 = x_1+x_2$，故D错误。故选：A。
\end{solution}
\end{question}

\begin{question}[points=4]
若集合 $\{(x,y) \mid y = x + t(x^2 - x), 0 \leq t \leq 1, 1 \leq x \leq 2\}$ 表示的图形中，两点间最大距离为 $d$、面积为 $S$，则（ ）
\begin{choices}
  \item $d = 3$，$S < 1$
  \item $d = 3$，$S > 1$
  \item $d = \sqrt{10}$，$S < 1$
  \item $d = \sqrt{10}$，$S > 1$
\end{choices}
\begin{solution}
\textbf{答案：}C

\textbf{分析：}先以 $t$ 为变量，分析可知所求集合表示的图形即为平面区域 $\begin{cases} y \leq x^2 \\ y \geq x \\ 1 \leq x \leq 2 \end{cases}$，结合图形分析求解即可。

\textbf{详解：}对任意给定 $x \in [1,2]$，则 $x^2 - x = x(x-1) \geq 0$，且 $t \in [0,1]$，可知 $x \leq x + t(x^2 - x) \leq x + (x^2 - x) = x^2$，即 $x \leq y \leq x^2$，再结合 $x$ 的任意性，所以所求集合表示的图形即为平面区域 $\begin{cases} y \leq x^2 \\ y \geq x \\ 1 \leq x \leq 2 \end{cases}$，如图阴影部分所示，其中 $A(1,1), B(2,2), C(2,4)$，可知任意两点间距离最大值 $d = AC = \sqrt{10}$；阴影部分面积 $S < S_{\triangle ABC} = \dfrac{1}{2} \times 1 \times 2 = 1$。故选：C。
\end{solution}
\end{question}

\section*{二、填空题}
本部分共5小题，每小题5分，共25分。

\begin{question}[points=5]
已知抛物线 $y^2 = 16x$，则焦点坐标为。
\fillin[{$(4,0)$}]
\begin{solution}
\textbf{答案：}$(4,0)$

\textbf{分析：}形如 $y^2 = 2px (p \neq 0)$ 的抛物线的焦点坐标为 $(\frac{p}{2},0)$，由此即可得解。

\textbf{详解：}由题意抛物线的标准方程为 $y^2 = 16x$，所以其焦点坐标为 $(4,0)$。
\end{solution}
\end{question}

\begin{question}[points=5]
已知 $\alpha \in \left[\dfrac{\pi}{6},\dfrac{\pi}{3}\right]$，且 $\alpha$ 与 $\beta$ 的终边关于原点对称，则 $\cos \beta$ 的最大值为。
\fillin[{$-\dfrac{1}{2}$}]
\begin{solution}
\textbf{答案：}$-\dfrac{1}{2}$

\textbf{分析：}首先得出 $\beta = \alpha + \pi + 2k\pi, k \in \mathbb{Z}$，结合三角函数单调性即可求解最值。

\textbf{详解：}由题意 $\beta = \alpha + \pi + 2k\pi, k \in \mathbb{Z}$，从而 $\cos \beta = \cos(\alpha+\pi+2k\pi) = -\cos \alpha$。因为 $\alpha \in [\dfrac{\pi}{6},\dfrac{\pi}{3}]$，所以 $\cos \alpha \in [\dfrac{1}{2},\dfrac{\sqrt{3}}{2}]$，$\cos \beta \in [-\dfrac{\sqrt{3}}{2},-\dfrac{1}{2}]$。当且仅当 $\alpha = \dfrac{\pi}{3}$，即 $\beta = \dfrac{4\pi}{3} + 2k\pi, k \in \mathbb{Z}$ 时，$\cos \beta$ 取得最大值，且最大值为 $-\dfrac{1}{2}$。
\end{solution}
\end{question}

\begin{question}[points=5]
已知双曲线 $\dfrac{x^2}{4} - y^2 = 1$，则过 $(3,0)$ 且和双曲线只有一个交点的直线的斜率为。
\fillin[{$\pm \dfrac{1}{2}$}]
\begin{solution}
\textbf{答案：}$\pm \dfrac{1}{2}$

\textbf{分析：}首先说明直线斜率存在，然后设出方程，联立双曲线方程，根据交点个数与方程根的情况列式即可求解。

\textbf{详解：}联立 $x=3$ 与 $\dfrac{x^2}{4} - y^2 = 1$，解得 $y = \pm \dfrac{\sqrt{5}}{2}$，这表明满足题意的直线斜率一定存在。设所求直线斜率为 $k$，则过点 $(3,0)$ 且斜率为 $k$ 的直线方程为 $y = k(x-3)$，联立 $\begin{cases} \dfrac{x^2}{4} - y^2 = 1 \\ y = k(x-3) \end{cases}$，化简并整理得：$(1-4k^2)x^2 + 24k^2x - 36k^2 - 4 = 0$。由题意得 $1-4k^2 = 0$ 或 $\Delta = (24k^2)^2 + 4(36k^2+4)(1-4k^2) = 0$，解得 $k = \pm \dfrac{1}{2}$ 或无解，即 $k = \pm \dfrac{1}{2}$，经检验，符合题意。
\end{solution}
\end{question}

\begin{question}[points=5]
已知三个圆柱的体积为公比为10的等比数列。第一个圆柱的直径为65mm，第二、三个圆柱的直径为325mm，第三个圆柱的高为230mm，求前两个圆柱的高度分别为。
\fillin[{$\dfrac{115}{2}\text{mm},23\text{mm}$}]
\begin{solution}
\textbf{答案：}$\dfrac{115}{2}\text{mm},23\text{mm}$

\textbf{分析：}根据体积为公比为10的等比数列可得关于高度的方程组，求出其解后可得前两个圆柱的高度。

\textbf{详解：}设第一个圆柱的高为 $h_1$，第二个圆柱的高为 $h_2$，则 $\dfrac{\pi\left(\dfrac{325}{2}\right)^2 h_2}{\pi\left(\dfrac{65}{2}\right)^2 h_1} = 10$，$\dfrac{\pi\left(\dfrac{325}{2}\right)^2 \times 230}{\pi\left(\dfrac{325}{2}\right)^2 h_2} = 10$，故 $h_2 = 23\text{mm}$，$h_1 = \dfrac{115}{2}\text{mm}$。
\end{solution}
\end{question}

\begin{question}[points=5]
已知 $M = \{k \mid a_k = b_k\}$，$a_n$，$b_n$ 不为常数列且各项均不相同，下列正确的是。

① $a_n$，$b_n$ 均为等差数列，则 $M$ 中最多一个元素；

② $a_n$，$b_n$ 均为等比数列，则 $M$ 中最多三个元素；

③ $a_n$ 为等差数列，$b_n$ 为等比数列，则 $M$ 中最多三个元素；

④ $a_n$ 单调递增，$b_n$ 单调递减，则 $M$ 中最多一个元素。
\fillin[{$①③④$}]
\begin{solution}
\textbf{答案：}$①③④$

\textbf{分析：}利用两类数列的散点图的特征可判断①④的正误，利用反例可判断②的正误，结合通项公式的特征及反证法可判断③的正误。

\textbf{详解：}对于①，因为 $\{a_n\}, \{b_n\}$ 均为等差数列，故它们的散点图分布在直线上，而两条直线至多有一个公共点，故 $M$ 中至多一个元素，故①正确。
对于②，取 $a_n = 2^{n-1}$，$b_n = -(-2)^{n-1}$，则 $\{a_n\}, \{b_n\}$ 均为等比数列，但当 $n$ 为偶数时，有 $a_n = 2^{n-1} = b_n = -(-2)^{n-1}$，此时 $M$ 中有无穷多个元素，故②错误。
对于③，设 $b_n = Aq^n$（$Aq \neq 0, q \neq \pm 1$），$a_n = kn + b$（$k \neq 0$）。若 $M$ 中至少四个元素，则关于 $n$ 的方程 $Aq^n = kn + b$ 至少有4个不同的正数解。若 $q>0, q\neq 1$，则由 $y=Aq^n$ 和 $y=kn+b$ 的散点图可得关于 $n$ 的方程 $Aq^n = kn+b$ 至多有两个不同的解，矛盾；若 $q<0, q\neq \pm 1$，考虑关于 $n$ 的方程 $Aq^n = kn+b$ 奇数解的个数和偶数解的个数。当 $Aq^n = kn+b$ 有偶数解，此方程即为 $A|q|^n = kn+b$，方程至多有两个偶数解，且有两个偶数解时 $Ak\ln|q|>0$，否则 $Ak\ln|q|<0$，因 $y=A|q|^n, y=kn+b$ 单调性相反，方程 $A|q|^n = kn+b$ 至多一个偶数解。当 $Aq^n = kn+b$ 有奇数解，此方程即为 $-A|q|^n = kn+b$，方程至多有两个奇数解，且有两个奇数解时 $-Ak\ln|q|>0$ 即 $Ak\ln|q|<0$，否则 $Ak\ln|q|>0$，因 $y=-A|q|^n, y=kn+b$ 单调性相反，方程 $-A|q|^n = kn+b$ 至多一个奇数解。因为 $Ak\ln|q|>0$，$Ak\ln|q|<0$ 不可能同时成立，故 $Aq^n = kn+b$ 不可能有4个不同的正数解，故③正确。
对于④，因为 $\{a_n\}$ 为单调递增，$\{b_n\}$ 为递减数列，前者散点图呈上升趋势，后者的散点图呈下降趋势，两者至多一个交点，故④正确。
故答案为：①③④。
\end{solution}
\end{question}

\section*{三、解答题}
本部分共6小题，共85分。解答应写出文字说明，演算步骤或证明过程。

\begin{question}[points=13]
在 $\triangle ABC$ 中，$a=7$，$A$ 为钝角，$\sin 2B = \dfrac{\sqrt{3}}{7} b \cos B$。

（1）求 $\angle A$；

（2）从条件①、条件②和条件③这三个条件中选择一个作为已知，求 $\triangle ABC$ 的面积。

① $b=7$；② $\cos B = \dfrac{13}{14}$；③ $c\sin A = \dfrac{5}{2}\sqrt{3}$。

注：如果选择条件①、条件②和条件③分别解答，按第一个解答计分。
\begin{solution}
\textbf{答案：}（1）$A=\dfrac{2\pi}{3}$；（2）选择①无解；选择②和③ $\triangle ABC$ 面积均为 $\dfrac{15\sqrt{3}}{4}$。

\textbf{分析：}（1）利用正弦定理即可求出答案；（2）选择①，利用正弦定理得 $B=\dfrac{\pi}{3}$，结合（1）问答案即可排除；选择②，首先求出 $\sin B = \dfrac{3\sqrt{3}}{14}$，再代入式子得 $b=3$，再利用两角和的正弦公式即可求出 $\sin C$，最后利用三角形面积公式即可；选择③，首先得到 $c=5$，再利用正弦定理得到 $\sin C = \dfrac{5\sqrt{3}}{14}$，再利用两角和的正弦公式即可求出 $\sin B$，最后利用三角形面积公式即可。

\textbf{详解：}（1）由题意得 $2\sin B \cos B = \dfrac{\sqrt{3}}{7} b \cos B$，因为 $A$ 为钝角，则 $\cos B \neq 0$，则 $2\sin B = \dfrac{\sqrt{3}}{7} b$，则 $\sin B = \dfrac{\sqrt{3}}{14} b$，由正弦定理 $\dfrac{b}{\sin B} = \dfrac{a}{\sin A}$ 得 $\dfrac{b}{\frac{\sqrt{3}}{14}b} = \dfrac{7}{\sin A}$，解得 $\sin A = \dfrac{\sqrt{3}}{2}$，因为 $A$ 为钝角，则 $A = \dfrac{2\pi}{3}$。
（2）选择① $b=7$，则 $\sin B = \dfrac{\sqrt{3}}{14} \times 7 = \dfrac{\sqrt{3}}{2}$，因为 $A=\dfrac{2\pi}{3}$，则 $B$ 为锐角，则 $B = \dfrac{\pi}{3}$，此时 $A+B = \pi$，不合题意，舍弃。
选择② $\cos B = \dfrac{13}{14}$，因为 $B$ 为三角形内角，则 $\sin B = \sqrt{1-\left(\dfrac{13}{14}\right)^2} = \dfrac{3\sqrt{3}}{14}$，则代入 $2\sin B = \dfrac{\sqrt{3}}{7} b$ 得 $2\times \dfrac{3\sqrt{3}}{14} = \dfrac{\sqrt{3}}{7} b$，解得 $b=3$。$\sin C = \sin(A+B) = \sin\left(\dfrac{2\pi}{3}+B\right) = \sin\dfrac{2\pi}{3}\cos B + \cos\dfrac{2\pi}{3}\sin B = \dfrac{\sqrt{3}}{2}\times\dfrac{13}{14} + \left(-\dfrac{1}{2}\right)\times\dfrac{3\sqrt{3}}{14} = \dfrac{5\sqrt{3}}{14}$。则 $S_{\triangle ABC} = \dfrac{1}{2}ab\sin C = \dfrac{1}{2}\times7\times3\times\dfrac{5\sqrt{3}}{14} = \dfrac{15\sqrt{3}}{4}$。
选择③ $c\sin A = \dfrac{5}{2}\sqrt{3}$，则有 $c\times\dfrac{\sqrt{3}}{2} = \dfrac{5}{2}\sqrt{3}$，解得 $c=5$。则由正弦定理得 $\dfrac{a}{\sin A} = \dfrac{c}{\sin C}$，即 $\dfrac{7}{\frac{\sqrt{3}}{2}} = \dfrac{5}{\sin C}$，解得 $\sin C = \dfrac{5\sqrt{3}}{14}$。因为 $C$ 为三角形内角，则 $\cos C = \sqrt{1-\left(\dfrac{5\sqrt{3}}{14}\right)^2} = \dfrac{11}{14}$。则 $\sin B = \sin(A+C) = \sin\left(\dfrac{2\pi}{3}+C\right) = \sin\dfrac{2\pi}{3}\cos C + \cos\dfrac{2\pi}{3}\sin C = \dfrac{\sqrt{3}}{2}\times\dfrac{11}{14} + \left(-\dfrac{1}{2}\right)\times\dfrac{5\sqrt{3}}{14} = \dfrac{3\sqrt{3}}{14}$。则 $S_{\triangle ABC} = \dfrac{1}{2}ac\sin B = \dfrac{1}{2}\times7\times5\times\dfrac{3\sqrt{3}}{14} = \dfrac{15\sqrt{3}}{4}$。
\end{solution}
\end{question}

\begin{question}[points=15]
已知四棱锥 $P-ABCD$，$AD\parallel BC$，$AB=BC=1$，$AD=3$，$DE=PE=2$，$E$ 是 $AD$ 上一点，$PE \perp AD$。

（1）若 $F$ 是 $PE$ 中点，证明：$BF\parallel$ 平面 $PCD$。

（2）若 $AB \perp$ 平面 $PED$，求平面 $PAB$ 与平面 $PCD$ 夹角的余弦值。

\begin{center}
  \includegraphics[width=.38\linewidth]{question-17-1.jpg}
\end{center}
\begin{solution}
\textbf{答案：}（1）证明见解析；（2）$\dfrac{\sqrt{30}}{30}$

\textbf{分析：}（1）取 $PD$ 的中点为 $S$，连接 $SF, SC$，可证四边形 $SFBC$ 为平行四边形，由线面平行的判定定理可得 $BF\parallel$ 平面 $PCD$。
（2）建立如图所示的空间直角坐标系，求出平面 $APB$ 和平面 $PCD$ 的法向量后可求夹角的余弦值。

\textbf{详解：}（1）取 $PD$ 的中点为 $S$，连接 $SF, SC$，则 $SF\parallel ED, SF = \dfrac{1}{2}ED = 1$，而 $ED\parallel BC, ED=2BC$，故 $SF\parallel BC, SF=BC$，故四边形 $SFBC$ 为平行四边形，故 $BF\parallel SC$，而 $BF \not\subset$ 平面 $PCD$，$SC \subset$ 平面 $PCD$，所以 $BF\parallel$ 平面 $PCD$。
（2）因为 $ED=2$，故 $AE=1$，故 $AE\parallel BC, AE=BC$，故四边形 $AECB$ 为平行四边形，故 $CE\parallel AB$，所以 $CE \perp$ 平面 $PAD$，而 $PE, ED \subset$ 平面 $PAD$，故 $CE \perp PE, CE \perp ED$，而 $PE \perp ED$，故建立如图所示的空间直角坐标系，则 $A(0,-1,0), B(1,-1,0), C(1,0,0), D(0,2,0), P(0,0,2)$，则 $\overrightarrow{PA}=(0,-1,-2), \overrightarrow{PB}=(1,-1,-2), \overrightarrow{PC}=(1,0,-2), \overrightarrow{PD}=(0,2,-2)$。设平面 $PAB$ 的法向量为 $\vec{m}=(x,y,z)$，则由 $\begin{cases} \vec{m}\cdot\overrightarrow{PA}=0 \\ \vec{m}\cdot\overrightarrow{PB}=0 \end{cases}$ 可得 $\begin{cases} -y-2z=0 \\ x-y-2z=0 \end{cases}$，取 $\vec{m}=(0,-2,1)$。设平面 $PCD$ 的法向量为 $\vec{n}=(a,b,c)$，则由 $\begin{cases} \vec{n}\cdot\overrightarrow{PC}=0 \\ \vec{n}\cdot\overrightarrow{PD}=0 \end{cases}$ 可得 $\begin{cases} a-2b=0 \\ 2b-2c=0 \end{cases}$，取 $\vec{n}=(2,1,1)$。故 $\cos\langle\vec{m},\vec{n}\rangle = \dfrac{-1}{\sqrt{5}\times\sqrt{6}} = -\dfrac{\sqrt{30}}{30}$，故平面 $PAB$ 与平面 $PCD$ 夹角的余弦值为 $\dfrac{\sqrt{30}}{30}$。
\end{solution}
\end{question}

\begin{question}[points=14]
已知某险种的保费为0.4万元，前3次出险每次赔付0.8万元，第4次赔付0.6万元。

赔偿次数：0，1，2，3，4

单数：800，100，60，30，10

在总体中抽样100单，以频率估计概率：

（1）求随机抽取一单，赔偿不少于2次的概率；

（2）（i）毛利润是保费与赔偿金额之差。设毛利润为 $X$，估计 $X$ 的数学期望；

（ii）若未赔偿过的保单下一保险期的保费下降 $4\%$，已赔偿过的增加 $20\%$。估计保单下一保险期毛利润的数学期望。
\begin{solution}
\textbf{答案：}（1）$\dfrac{1}{10}$；（2）(i)0.122万元；(ii)0.1252万元

\textbf{分析：}（1）根据题设中的数据可求赔偿次数不少于2的概率；（2）（i）设 $\xi$ 为赔付金额，则 $\xi$ 可取0,0.8,1.6,2.4,3，用频率估计概率后可求 $\xi$ 的分布列及数学期望，从而可求 $E(X)$；（ii）先算出下一期保费的变化情况，结合（1）的结果可求 $E(Y)$。

\textbf{详解：}（1）设 $A$ 为“随机抽取一单，赔偿不少于2次”，由题设中的统计数据可得 $P(A) = \dfrac{60+30+10}{800+100+60+30+10} = \dfrac{1}{10}$。
（2）（i）设 $\xi$ 为赔付金额，则 $\xi$ 可取0,0.8,1.6,2.4,3，由题设中的统计数据可得 $P(\xi=0)=\dfrac{800}{1000}=\dfrac{4}{5}$，$P(\xi=0.8)=\dfrac{100}{1000}=\dfrac{1}{10}$，$P(\xi=1.6)=\dfrac{60}{1000}=\dfrac{3}{50}$，$P(\xi=2.4)=\dfrac{30}{1000}=\dfrac{3}{100}$，$P(\xi=3)=\dfrac{10}{1000}=\dfrac{1}{100}$。故 $E(\xi)=0\times\dfrac{4}{5}+0.8\times\dfrac{1}{10}+1.6\times\dfrac{3}{50}+2.4\times\dfrac{3}{100}+3\times\dfrac{1}{100}=0.278$。故 $E(X)=0.4-0.278=0.122$（万元）。
（ii）由题设保费的变化为 $0.4\times\dfrac{4}{5}\times96\% + 0.4\times\dfrac{1}{5}\times1.2 = 0.4032$，故 $E(Y)=0.122+0.4032-0.4=0.1252$（万元）。
\end{solution}
\end{question}

\begin{question}[points=14]
已知椭圆方程 $C:\dfrac{x^2}{a^2}+\dfrac{y^2}{b^2}=1 (a>b>0)$，焦点和短轴端点构成边长为2的正方形，过 $(0,t)(t>\sqrt{2})$ 的直线 $l$ 与椭圆交于 $A,B$，$C(0,1)$，连接 $AC$ 交椭圆于 $D$。

（1）求椭圆方程和离心率；

（2）若直线 $BD$ 的斜率为0，求 $t$。
\begin{solution}
\textbf{答案：}（1）$\dfrac{x^2}{4}+\dfrac{y^2}{2}=1, e=\dfrac{\sqrt{2}}{2}$；（2）$t=2$

\textbf{分析：}（1）由题意得 $b=c=\sqrt{2}$，进一步得 $a$，由此即可得解；（2）说明直线 $AB$ 斜率存在，设 $AB:y=kx+t, t>\sqrt{2}$，$A(x_1,y_1),B(x_2,y_2)$，联立椭圆方程，由韦达定理有 $x_1+x_2=\dfrac{-4kt}{1+2k^2}, x_1x_2=\dfrac{2t^2-4}{2k^2+1}$，而 $AD: y=\dfrac{y_1-y_2}{x_1+x_2}(x-x_1)+y_1$，令 $x=0$，即可得解。

\textbf{详解：}（1）由题意 $b=c=\sqrt{\dfrac{2^2}{2}}=\sqrt{2}$，从而 $a=\sqrt{b^2+c^2}=2$，所以椭圆方程为 $\dfrac{x^2}{4}+\dfrac{y^2}{2}=1$，离心率为 $e=\dfrac{\sqrt{2}}{2}$。
（2）显然直线 $AB$ 斜率存在，否则 $B,D$ 重合，直线 $BD$ 斜率不存在与题意不符，同样直线 $AB$ 斜率不为0，否则直线 $AB$ 与椭圆无交点，矛盾。从而设 $AB: y=kx+t, t>\sqrt{2}$，$A(x_1,y_1),B(x_2,y_2)$。联立 $\begin{cases} \dfrac{x^2}{4}+\dfrac{y^2}{2}=1 \\ y=kx+t \end{cases}$，化简并整理得 $(1+2k^2)x^2+4ktx+2t^2-4=0$，由题意 $\Delta=16k^2t^2-8(2k^2+1)(t^2-2)=8(4k^2+2-t^2)>0$，即 $k,t$ 应满足 $4k^2+2-t^2>0$。所以 $x_1+x_2=\dfrac{-4kt}{1+2k^2}$，$x_1x_2=\dfrac{2t^2-4}{2k^2+1}$。若直线 $BD$ 斜率为0，由椭圆的对称性可设 $D(-x_2,y_2)$，所以 $AD: y=\dfrac{y_1-y_2}{x_1+x_2}(x-x_1)+y_1$，在直线 $AD$ 方程中令 $x=0$，得 $y_C=\dfrac{x_1y_2+x_2y_1}{x_1+x_2}=\dfrac{x_1(kx_2+t)+x_2(kx_1+t)}{x_1+x_2}=\dfrac{2kx_1x_2+t(x_1+x_2)}{x_1+x_2}=\dfrac{2k\cdot\frac{2t^2-4}{2k^2+1}+t\cdot\frac{-4kt}{1+2k^2}}{\frac{-4kt}{1+2k^2}} = \dfrac{\frac{4k(t^2-2)-4k t^2}{1+2k^2}}{\frac{-4kt}{1+2k^2}} = \dfrac{-8k}{(1+2k^2)}\cdot\dfrac{1+2k^2}{-4kt} = \dfrac{2}{t} = 1$，所以 $t=2$。此时 $k$ 应满足 $\begin{cases} 4k^2+2-t^2=4k^2-2>0 \\ k\neq0 \end{cases}$，即 $k< -\dfrac{\sqrt{2}}{2}$ 或 $k>\dfrac{\sqrt{2}}{2}$。综上所述，$t=2$ 满足题意。
\end{solution}
\end{question}

\begin{question}[points=15]
已知 $f(x)=x+k\ln(1+x)$ 在 $(t,f(t))(t>0)$ 处切线为 $l$。

（1）若切线 $l$ 的斜率 $k=-1$，求 $f(x)$ 单调区间；

（2）证明：切线 $l$ 不经过 $(0,0)$；

（3）已知 $k=1$，$A(t,f(t))$，$C(0,f(t))$，$O(0,0)$，其中 $t>0$，切线 $l$ 与 $y$ 轴交于点 $B$ 时。当 $2S_{\triangle ACO}=15S_{\triangle ABO}$，符合条件的 $A$ 的个数为？

（参考数据：$1.09<\ln3<1.10$，$1.60<\ln5<1.61$，$1.94<\ln7<1.95$）
\begin{solution}
\textbf{答案：}（1）单调递减区间为 $(-1,0)$，单调递增区间为 $(0,+\infty)$；（2）证明见解析；（3）2

\textbf{分析：}（1）直接代入 $k=-1$，再利用导数研究其单调性即可；（2）写出切线方程 $y-f(t)=\left(1+\dfrac{k}{1+t}\right)(x-t)(t>0)$，将 $(0,0)$ 代入再设新函数 $F(t)=\ln(1+t)-\dfrac{t}{1+t}$，利用导数研究其零点即可；（3）分别写出面积表达式，代入 $2S_{\triangle ACO}=15S_{\triangle ABO}$ 得到 $13\ln(1+t)-2t-15\dfrac{t}{1+t}=0$，再设新函数 $h(t)=13\ln(1+t)-2t-\dfrac{15t}{1+t}(t>0)$ 研究其零点即可。

\textbf{详解：}（1）当 $k=-1$ 时，$f(x)=x-\ln(1+x)$，$f'(x)=1-\dfrac{1}{1+x}=\dfrac{x}{1+x}(x>-1)$。当 $x\in(-1,0)$ 时，$f'(x)<0$；当 $x\in(0,+\infty)$ 时，$f'(x)>0$。所以 $f(x)$ 的单调递减区间为 $(-1,0)$，单调递增区间为 $(0,+\infty)$。
（2）$f'(x)=1+\dfrac{k}{1+x}$，切线 $l$ 的斜率为 $1+\dfrac{k}{1+t}$，则切线方程为 $y-f(t)=\left(1+\dfrac{k}{1+t}\right)(x-t)(t>0)$。将 $(0,0)$ 代入则 $-f(t)=-t\left(1+\dfrac{k}{1+t}\right)$，$f(t)=t\left(1+\dfrac{k}{1+t}\right)$，即 $t+k\ln(1+t)=t+\dfrac{kt}{1+t}$，则 $\ln(1+t)=\dfrac{t}{1+t}$，$\ln(1+t)-\dfrac{t}{1+t}=0$。令 $F(t)=\ln(1+t)-\dfrac{t}{1+t}$，假设 $l$ 过 $(0,0)$，则 $F(t)$ 在 $t>0$ 存在零点。$F'(t)=\dfrac{1}{1+t}-\dfrac{(1+t)-t}{(1+t)^2}=\dfrac{t}{(1+t)^2}>0$，所以 $F(t)$ 在 $(0,+\infty)$ 上单调递增，$F(t)>F(0)=0$，$F(t)$ 在 $(0,+\infty)$ 无零点，与假设矛盾，故直线 $l$ 不过 $(0,0)$。
（3）$k=1$ 时，$f(x)=x+\ln(1+x)$，$f'(x)=1+\dfrac{1}{1+x}=\dfrac{x+2}{1+x}>0$。$S_{\triangle ACO}=\dfrac{1}{2}tf(t)$，设 $l$ 与 $y$ 轴交点 $B$ 为 $(0,q)$，$t>0$ 时，若 $q<0$，则此时 $l$ 与 $f(x)$ 必有交点，与切线定义矛盾。由（2）知 $q\neq0$，所以 $q>0$。则切线 $l$ 的方程为 $y-t-\ln(t+1)=\left(1+\dfrac{1}{1+t}\right)(x-t)$，令 $x=0$，则 $y=q=\ln(1+t)-\dfrac{t}{t+1}$。因为 $2S_{\triangle ACO}=15S_{\triangle ABO}$，则 $2tf(t)=15t\left[\ln(1+t)-\dfrac{t}{t+1}\right]$，即 $13\ln(1+t)-2t-15\dfrac{t}{1+t}=0$。记 $h(t)=13\ln(1+t)-2t-\dfrac{15t}{1+t}(t>0)$，则满足条件的 $A$ 有几个即 $h(t)$ 有几个零点。$h'(t)=\dfrac{13}{1+t}-2-\dfrac{15}{(1+t)^2}=\dfrac{13(1+t)-2(1+t)^2-15}{(1+t)^2}=\dfrac{-2t^2+9t-4}{(1+t)^2}=\dfrac{(-2t+1)(t-4)}{(1+t)^2}$。当 $t\in(0,\frac{1}{2})$ 时，$h'(t)<0$，$h(t)$ 单调递减；当 $t\in(\frac{1}{2},4)$ 时，$h'(t)>0$，$h(t)$ 单调递增；当 $t\in(4,+\infty)$ 时，$h'(t)<0$，$h(t)$ 单调递减。因为 $h(0)=0$，$h(\frac{1}{2})<0$，$h(4)=13\ln5-20>13\times1.6-20=0.8>0$，$h(24)=13\ln25-48-\dfrac{15\times24}{25}=26\ln5-48-\dfrac{72}{5}<26\times1.61-48-14.4=41.86-62.4=-20.54<0$，所以由零点存在性定理及 $h(t)$ 的单调性，$h(t)$ 在 $(\frac{1}{2},4)$ 上必有一个零点，在 $(4,24)$ 上必有一个零点。综上所述，$h(t)$ 有两个零点，即满足 $2S_{\triangle ACO}=15S_{\triangle ABO}$ 的 $A$ 有两个。
\end{solution}
\end{question}

\begin{question}[points=14]
设集合 $M = \{(i,j,s,t) \mid i\in\{1,2\}, j\in\{3,4\}, s\in\{5,6\}, t\in\{7,8\}, 2|(i+j+s+t)\}$。对于给定有穷数列 $A:\{a_n\}(1\leq n\leq 8)$，及序列 $\Omega:\omega_1,\omega_2,\dots,\omega_s$，$\omega_k=(i_k,j_k,s_k,t_k)\in M$，定义变换 $T$：将数列 $A$ 的第 $i_1,j_1,s_1,t_1$ 项加1，得到数列 $T_1(A)$；将数列 $T_1(A)$ 的第 $i_2,j_2,s_2,t_2$ 项加1，得到数列 $T_2T_1(A)$；重复上述操作，得到数列 $T_s\cdots T_2T_1(A)$，记为 $\Omega(A)$。

（1）给定数列 $A:1,3,2,4,6,3,1,9$ 和序列 $\Omega:(1,3,5,7),(2,4,6,8),(1,3,5,7)$，写出 $\Omega(A)$；

（2）是否存在序列 $\Omega$，使得 $\Omega(A)$ 为 $a_1+2, a_2+6, a_3+4, a_4+2, a_5+8, a_6+2, a_7+4, a_8+4$，若存在，写出一个符合条件的 $\Omega$；若不存在，请说明理由；

（3）若数列 $A$ 的各项均为正整数，且 $a_1+a_3+a_5+a_7$ 为偶数，证明：“存在序列 $\Omega$，使得 $\Omega(A)$ 为常数列”的充要条件为“$a_1+a_2=a_3+a_4=a_5+a_6=a_7+a_8$”。
\begin{solution}
\textbf{答案：}（1）$\Omega(A):3,4,4,5,8,4,3,10$；（2）不存在，理由见解析；（3）证明见解析

\textbf{分析：}（1）直接按照 $\Omega(A)$ 的定义写出 $\Omega(A)$ 即可；（2）利用反证法，假设存在符合条件的 $\Omega$，由此列出方程组，进一步说明方程组无解即可；（3）分充分性和必要性两方面论证。

\textbf{详解：}（1）$\Omega(A):3,4,4,5,8,4,3,10$。
（2）假设存在符合条件的 $\Omega$，可知 $\Omega(A)$ 的第1,2项之和为 $a_1+a_2+s$，第3,4项之和为 $a_3+a_4+s$，则 $\begin{cases} (a_1+2)+(a_2+6)=a_1+a_2+s \ (a_3+4)+(a_4+2)=a_3+a_4+s \end{cases}$，而该方程组无解，故假设不成立，故不存在符合条件的 $\Omega$。
（3）设序列 $T_k\cdots T_2T_1(A)$ 为 $\{a_{k,n}\}(1\leq n\leq 8)$，特别规定 $a_{0,n}=a_n(1\leq n\leq 8)$。
必要性：若存在序列 $\Omega:\omega_1,\omega_2,\dots,\omega_s$，使得 $\Omega(A)$ 为常数列，则 $a_{s,1}=a_{s,2}=a_{s,3}=a_{s,4}=a_{s,5}=a_{s,6}=a_{s,7}=a_{s,8}$，所以 $a_{s,1}+a_{s,2}=a_{s,3}+a_{s,4}=a_{s,5}+a_{s,6}=a_{s,7}+a_{s,8}$。根据 $T_k\cdots T_2T_1(A)$ 的定义，显然有 $a_{k,2j-1}+a_{k,2j}=a_{k-1,2j-1}+a_{k-1,2j}$，这里 $j=1,2,3,4$，$k=1,2,\dots$。所以不断使用该式就得到 $a_1+a_2=a_3+a_4=a_5+a_6=a_7+a_8$，必要性得证。
充分性：若 $a_1+a_2=a_3+a_4=a_5+a_6=a_7+a_8$。由已知 $a_1+a_3+a_5+a_7$ 为偶数，而 $a_1+a_2=a_3+a_4=a_5+a_6=a_7+a_8$，所以 $a_2+a_4+a_6+a_8=4(a_1+a_2)-(a_1+a_3+a_5+a_7)$ 也是偶数。我们设 $T_s\cdots T_2T_1(A)$ 是通过合法的序列 $\Omega$ 的变换能得到的所有可能的数列 $\Omega(A)$ 中，使得 $a_{s,1}-a_{s,2}+a_{s,3}-a_{s,4}+a_{s,5}-a_{s,6}+a_{s,7}-a_{s,8}$ 最小的一个。上面已经证明 $a_{k,2j-1}+a_{k,2j}=a_{k-1,2j-1}+a_{k-1,2j}$，这里 $j=1,2,3,4$，$k=1,2,\dots$。从而由 $a_1+a_2=a_3+a_4=a_5+a_6=a_7+a_8$ 可得 $a_{s,1}+a_{s,2}=a_{s,3}+a_{s,4}=a_{s,5}+a_{s,6}=a_{s,7}+a_{s,8}$。同时，由于 $i_k+j_k+s_k+t_k$ 总是偶数，所以 $a_{k,1}+a_{k,3}+a_{k,5}+a_{k,7}$ 和 $a_{k,2}+a_{k,4}+a_{k,6}+a_{k,8}$ 的奇偶性保持不变，从而 $a_{s,1}+a_{s,3}+a_{s,5}+a_{s,7}$ 和 $a_{s,2}+a_{s,4}+a_{s,6}+a_{s,8}$ 都是偶数。
下面证明不存在 $j=1,2,3,4$ 使得 $a_{s,2j-1}-a_{s,2j}\geq 2$。假设存在，根据对称性，不妨设 $j=1$，$a_{s,2j-1}-a_{s,2j}\geq 2$，即 $a_{s,1}-a_{s,2}\geq 2$。
情况1：若 $a_{s,3}-a_{s,4}+a_{s,5}-a_{s,6}+a_{s,7}-a_{s,8}=0$，则由 $a_{s,1}+a_{s,3}+a_{s,5}+a_{s,7}$ 和 $a_{s,2}+a_{s,4}+a_{s,6}+a_{s,8}$ 都是偶数，知 $a_{s,1}-a_{s,2}\geq 4$。对该数列连续作四次变换 $(2,3,5,8)$，$(2,4,6,8)$，$(2,3,6,7)$，$(2,4,5,7)$ 后，新的 $a_{s+4,1}-a_{s+4,2}+a_{s+4,3}-a_{s+4,4}+a_{s+4,5}-a_{s+4,6}+a_{s+4,7}-a_{s+4,8}$ 相比原来的 $a_{s,1}-a_{s,2}+a_{s,3}-a_{s,4}+a_{s,5}-a_{s,6}+a_{s,7}-a_{s,8}$ 减少4，这与 $a_{s,1}-a_{s,2}+a_{s,3}-a_{s,4}+a_{s,5}-a_{s,6}+a_{s,7}-a_{s,8}$ 的最小性矛盾。
情况2：若 $a_{s,3}-a_{s,4}+a_{s,5}-a_{s,6}+a_{s,7}-a_{s,8}>0$，不妨设 $a_{s,3}-a_{s,4}>0$。
情况2-1：如果 $a_{s,3}-a_{s,4}\geq 1$，则对该数列连续作两次变换 $(2,4,5,7)$，$(2,4,6,8)$ 后，新的 $a_{s+2,1}-a_{s+2,2}+a_{s+2,3}-a_{s+2,4}+a_{s+2,5}-a_{s+2,6}+a_{s+2,7}-a_{s+2,8}$ 相比原来的 $a_{s,1}-a_{s,2}+a_{s,3}-a_{s,4}+a_{s,5}-a_{s,6}+a_{s,7}-a_{s,8}$ 至少减少2，这与最小性矛盾。
情况2-2：如果 $a_{s,4}-a_{s,3}\geq 1$，则对该数列连续作两次变换 $(2,3,5,8)$，$(2,3,6,7)$ 后，新的 $a_{s+2,1}-a_{s+2,2}+a_{s+2,3}-a_{s+2,4}+a_{s+2,5}-a_{s+2,6}+a_{s+2,7}-a_{s+2,8}$ 相比原来的至少减少2，这与最小性矛盾。
这就说明无论如何都会导致矛盾，所以对任意的 $j=1,2,3,4$ 都有 $a_{s,2j-1}-a_{s,2j}\leq 1$。
假设存在 $j=1,2,3,4$ 使得 $a_{s,2j-1}-a_{s,2j}=1$，则 $a_{s,2j-1}+a_{s,2j}$ 是奇数，所以 $a_{s,1}+a_{s,2}=a_{s,3}+a_{s,4}=a_{s,5}+a_{s,6}=a_{s,7}+a_{s,8}$ 都是奇数，设为 $2N+1$。则此时对任意 $j=1,2,3,4$，由 $|a_{s,2j-1}-a_{s,2j}|\leq 1$ 可知必有 $\{a_{s,2j-1},a_{s,2j}\}=\{N,N+1\}$。而 $a_{s,1}+a_{s,3}+a_{s,5}+a_{s,7}$ 和 $a_{s,2}+a_{s,4}+a_{s,6}+a_{s,8}$ 都是偶数，故集合 $\{m|a_{s,m}=N\}$ 中的四个元素 $i,j,s,t$ 之和为偶数，对该数列进行一次变换 $(i,j,s,t)$，则该数列成为常数列，新的 $a_{s+1,1}-a_{s+1,2}+a_{s+1,3}-a_{s+1,4}+a_{s+1,5}-a_{s+1,6}+a_{s+1,7}-a_{s+1,8}$ 等于零，比原来的 $a_{s,1}-a_{s,2}+a_{s,3}-a_{s,4}+a_{s,5}-a_{s,6}+a_{s,7}-a_{s,8}$ 更小，这与最小性矛盾。
综上，只可能 $a_{s,2j-1}-a_{s,2j}=0$（$j=1,2,3,4$），而 $a_{s,1}+a_{s,2}=a_{s,3}+a_{s,4}=a_{s,5}+a_{s,6}=a_{s,7}+a_{s,8}$，故 $\{a_{s,n}\}=\Omega(A)$ 是常数列，充分性得证。
\end{solution}
\end{question}

